\documentclass[11pt]{article}
\usepackage[letterpaper,margin=0.9in]{geometry}
\usepackage{amsmath,amssymb,amsthm,mathtools,booktabs,microtype}
\usepackage[T1]{fontenc}
\usepackage{lmodern}
\usepackage[hidelinks]{hyperref}
\hypersetup{pdftitle={One Prime Signal for All Dirichlet Zeros},pdfsubject={Universal multiplicity reconstruction from smoothed autocorrelation}}
\usepackage{enumitem}
\allowdisplaybreaks[2]
\setlength{\emergencystretch}{2em}
\newtheorem{theorem}{Theorem}
\newtheorem{lemma}[theorem]{Lemma}
\newtheorem{proposition}[theorem]{Proposition}
\newtheorem{corollary}[theorem]{Corollary}
\theoremstyle{remark}\newtheorem{remark}[theorem]{Remark}
\newcommand{\R}{\mathbb R}\newcommand{\C}{\mathbb C}\newcommand{\N}{\mathbb N}\newcommand{\Z}{\mathbb Z}\newcommand{\Q}{\mathbb Q}
\newcommand{\ord}{\operatorname{ord}}\newcommand{\Rea}{\operatorname{Re}}\newcommand{\e}{\mathrm e}
\title{One Prime Signal for All Dirichlet Zeros\\\large Universal multiplicity reconstruction from a smoothed autocorrelation}
\author{Research note}
\date{26 November 2026}
\begin{document}
\maketitle
\begin{abstract}
Assuming the zero-location theorem for Dirichlet $L$-functions established in the Alp\"oge--Claude Lean development, we construct one complex-valued weight on the positive integers, with uniformly positive real part, such that the autocorrelation of its normalized smoothed prime-power counting function determines the entire character-labelled nontrivial zero multiset of every primitive Dirichlet $L$-function. The same result holds with primes only, after an explicit constant correction for prime squares. No simplicity, noncoincidence, spacing, or linear independence hypothesis on zeros is needed. The main new ingredient is a countable, positive algebraic code that separates multiplicity sequences of at most conductor-linear growth. The analytic argument gives an absolutely summable Fourier expansion, even across all conductors; Bohr averaging recovers its power at each frequency, from which the code is decoded. We prove a prefix stability bound and explain the distinction between exact reconstruction and practical numerical recovery. This is a proposed original theorem, not a claim of established bibliographic priority.
\end{abstract}

\section{Source analysis and scope}
The two requested sources are the repository \texttt{lsernau/grh-lean} \cite{repo} and the Riemann hypothesis survey \cite{wiki}. The repository was inspected through GitHub at commit \texttt{72fead1} (full hash in Reference~[1]). Its statements and final exports establish RH and GRH for Dirichlet $L$-functions (all positive moduli and all characters). In particular, \path{LFunction_ne_zero_of_half_lt_re} excludes zeros in the open right half-plane $\Rea s>1/2$, except that the principal-character pole at $s=1$ is excluded from the statement. The primitive-character critical-strip theorem is also explicitly exported. This is not a theorem for arbitrary Hecke or automorphic $L$-functions.

The readable edition's guide and the files \path{New/Reduction/MainInequality.lean} and \path{New/Reduction/OneSidedBound.lean} expose an especially useful interface. The former compares an index-corrected covolume from above and below; size and divisibility gains cancel a common main term. The latter obtains a one-sided estimate for a remainder based on
\[
\sum_{\substack{n\le x\\ n\equiv1\ (q)}}\Lambda(n)(1-\chi_K(n))\left(\frac1n-\frac1x\right)
\]
after subtracting its logarithmic main term and constant. A zero-detection argument then gives half-plane nonvanishing. We take two ideas from this structure: smoothing prime powers to obtain convergent analytic transforms, and keeping weights at the relevant spectral objects positive to prevent cancellation. Our smoothing kernel is different, namely $(1-n/x)_+$, and our positivity is on character coefficients, not the original lattice construction. We use the final GRH theorem as an input, not the unexported arithmetic geometry as a black box for a different conjecture.

The survey \cite{wiki} emphasizes the critical-line description, the explicit formula linking zeros to prime-counting oscillations, and classical analytic criteria such as Nyman--Beurling. The present question is inverse: how much zero information is retained if only second-order statistics of one prime signal are kept? Existing explicit-formula and almost-periodic work, including Akbary--Ng--Shahabi \cite{ans} and Rubinstein--Sarnak \cite{rs}, gives the natural background. The contribution here is the countable noncancelling encoding, which permits recovering labelled multiplicities even at coincident frequencies. A targeted literature search did not locate this precise universal reconstruction statement; the ingredients are classical and priority has not been certified.

The repository reports machine checks with the three standard Lean axioms, and independent checks of the original and shortened developments are documented by Rubinstein \cite{rub}. We inspected the statements, guide, and reduction files, but did not rebuild the repository or formally verify this extension. Throughout the paper GRH is an explicit named input; the proof below is a classical mathematical derivation from that input. No claim of solving a further previously named open conjecture is made.

\section{The theorem}
Enumerate the primitive Dirichlet characters as $(\chi_j)_{j\ge0}$, first by conductor, then in any fixed effective order within a conductor. Write $q_j$ for the conductor and $w_j=q_j+1$. The character $\chi_0=1$ has conductor $1$ and $L(s,\chi_0)=\zeta(s)$. Extend all characters to integers in the usual way, with zero on nonunits. For $\gamma\in\R$ write
\[
m_j(\gamma)=\ord_{s=1/2+i\gamma}L(s,\chi_j),
\]
where the order is zero when the function is nonzero. All these orders are nonnegative integers.

\begin{theorem}[Universal zero reconstruction]\label{thm:main}
There exist explicitly specified positive real algebraic coefficients $(c_j)_{j\ge0}$, with $c_0=1$ and $\sum_{j\ge1}c_j w_j\le1/4$, such that the following holds under Dirichlet GRH. Define
\begin{align*}
W(n)&=\sum_{j\ge0}c_j\chi_j(n),\\
S(x)&=\sum_{n\le x}\Lambda(n)W(n)(1-n/x),\\
F(t)&=\e^{-t/2}\bigl(S(\e^t)-\e^t/2\bigr).
\end{align*}
Here $\Lambda$ is the von Mangoldt function. In particular $\Rea W(n)\ge3/4$ and $|W(n)|\le5/4$. The limits
\begin{align}
C(h)&=\lim_{T\to\infty}\frac1T\int_0^T F(t+h)\overline{F(t)}\,dt,\qquad h\in\R,\label{eq:c}\\
P(\gamma)&=\lim_{U\to\infty}\frac1U\int_0^U C(h)\e^{-i\gamma h}\,dh,\qquad\gamma\in\R,\label{eq:p}
\end{align}
exist. With $\rho_\gamma=1/2+i\gamma$, they satisfy
\begin{equation}\label{eq:recover}
\boxed{\quad |\rho_\gamma(\rho_\gamma+1)|\sqrt{P(\gamma)}
=\sum_{j\ge0}c_j m_j(\gamma).\quad}
\end{equation}
The right-hand side uniquely determines the entire sequence $(m_j(\gamma))_{j\ge0}$. Thus $C$ alone, together with the fixed public code, determines the labelled nontrivial zeros of all primitive Dirichlet $L$-functions, including multiplicity. This also determines the critical-strip zeros of every imprimitive Dirichlet $L$-function through its inducing primitive character.
\end{theorem}

\begin{remark}
The signal is scalar and complex-valued; equivalently it consists of two real channels. The conjugate in \eqref{eq:c} is essential, and the full complex correlation, rather than just its even real part or its value at zero lag, is retained. All sums defining $S(x)$ over integers are finite, and $W(n)$ is uniformly computably approximable by finite character sums. ``One signal'' does not mean one finite observation or one unweighted arithmetic progression.
\end{remark}

\section{A countable positive code}\label{sec:code}
We first prove a purely algebraic lemma, independent of GRH. For a sequence of positive integer weights $(w_j)$ with $w_0=2$, put
\[
\mathcal M=\bigcup_{H=1}^\infty \prod_{j=0}^\infty \{0,1,\ldots,H w_j\}.
\]
Let $p_j$ be the $j$th rational prime for $j\ge1$. We construct $c_0=1$, $c_j=2^{-k_j}\sqrt{p_j}$ with nonnegative integers $k_j$ and define $k_0=0$. Once the first $n+1$ coefficients are defined, put
\begin{equation}\label{eq:delta}
D_n=2^{\max_{0\le j\le n} k_j},\qquad
 d_n=\frac1{D_n\bigl(3 D_n(n+1)\bigr)^{2^n-1}}.
\end{equation}
For $j\ge1$, choose the least nonnegative $k_j$ such that
\begin{equation}\label{eq:rule}
2^{-k_j}p_j w_j\le 2^{-j-2}
\min\left(1,\min_{0\le n<j}\frac{d_n}{n+1}\right).
\end{equation}
This is an entirely rational finite computation at each step, and a suitable integer always exists. It implies
\begin{align}
\sum_{j\ge1}c_jw_j&\le1/4,\label{eq:total}\\
\sum_{j>n}c_jw_j&\le\frac{d_n}{n+1}2^{-n-2}.
\label{eq:tail}
\end{align}
Indeed $\sqrt{p_j}\le p_j$ and the geometric series in \eqref{eq:rule} sums to $2^{-n-2}$ from $j=n+1$ onward.

\begin{lemma}[Prefix separation]\label{lem:sep}
If integers $u_0,\ldots,u_n$ are not all zero and $|u_j|\le(n+1)w_j$, then
\[
\left|\sum_{j=0}^n c_j u_j\right|\ge d_n.
\]
\end{lemma}
\begin{proof}
The numbers $1,\sqrt{p_1},\ldots,\sqrt{p_n}$ are linearly independent over $\Q$. In fact the field $K_n=\Q(\sqrt{p_1},\ldots,\sqrt{p_n})$ has degree $2^n$, with independent sign changes as automorphisms. For completeness, induct on $n$. If $\sqrt{p_n}$ belonged to $K_{n-1}$, it would be a simultaneous eigenvector of the sign changes, since its square is rational. The products of subsets of the preceding square roots form an eigenbasis with distinct sign characters, so it would be a rational multiple of one such product. Squaring contradicts the parity of the $p_n$-valuation. The extension therefore doubles the degree, and the independent sign changes extend. This proves both claims.

The number $z=D_n\sum_{j=0}^n c_j u_j$ is a nonzero algebraic integer of $K_n$. By \eqref{eq:total}, each conjugate satisfies
\[
|\sigma(z)|\le D_n(n+1)\sum_{j=0}^n c_jw_j
\le\tfrac94 D_n(n+1)<3 D_n(n+1).
\]
Its field norm is a nonzero rational integer. Bounding all conjugates except the specified real one gives $1\le|z|(3 D_n(n+1))^{2^n-1}$. Division by $D_n$ proves the stated (slightly conservative) bound.
\end{proof}

\begin{proposition}[Injectivity and stable prefix decoding]\label{prop:code}
The map $E:\mathcal M\to\R_{\ge0}$ defined by $E(m)=\sum_j c_j m_j$ is injective. If $m_j\le H w_j$ and $n+1\ge H$, its first $n+1$ entries are uniquely identifiable from any real $A'$ with
\[
|A'-E(m)|<d_n/8.
\]
The decoder searches all prefixes $v$ in $0\le v_j\le(n+1)w_j$ and chooses the unique one satisfying $|A'-\sum_{j=0}^n c_j v_j|<d_n/2$.
\end{proposition}
\begin{proof}
Absolute convergence follows from \eqref{eq:total}. The true prefix differs from $E(m)$ by at most $H\sum_{j>n}c_j w_j\le d_n/4$, by \eqref{eq:tail}; hence it lies within $3d_n/8$ of $A'$. By Lemma~\ref{lem:sep}, any different prefix is at least $d_n$ away from the true prefix, so cannot be within $d_n/2$ of $A'$. For two distinct sequences in $\mathcal M$, choose $H$ bounding both and $n\ge H-1$ at least as large as their first differing index. Their prefixes are different, yet equal encoded values would force the same decoded prefix. This contradiction proves injectivity. The finite search can use certified rational intervals for the square roots, since the admissible prefix is strictly separated from the decision boundary.
\end{proof}

\section{A uniform smoothed explicit formula}
The following details are included to make the passage from finitely many characters to all conductors rigorous. We use the standard functional equation, Euler product, entire order-one completed functions, Hadamard factorization and Stirling estimates (see, e.g., \cite{mv,dav}). All implied constants in the next lemma are absolute. No zero-density theorem or effective spacing hypothesis is required.

\begin{lemma}\label{lem:explicit}
Assume GRH for a primitive character $\chi$ of conductor $q$ (including the conductor-one character). Let $\delta_\chi=1$ if $q=1$ and $0$ otherwise. For $x\ge1$,
\begin{equation}\label{eq:explicit}
\sum_{n\le x}\Lambda(n)\chi(n)(1-n/x)
=\delta_\chi\frac{x}{2}-\sum_{\rho}\frac{x^\rho}{\rho(\rho+1)}
+O\bigl(x^{1/4}\log(q+2)\bigr),
\end{equation}
where the sum is over all nontrivial zeros with multiplicity. Moreover for $Y\ge1$,
\begin{equation}\label{eq:zerotail}
\sum_{|\Im\rho|>Y}\frac1{|\rho(\rho+1)|}
\ll\frac{\log(q(Y+2))}{Y+1},\qquad
\sum_\rho\frac1{|\rho(\rho+1)|}\ll\log(q+2).
\end{equation}
Each interval $[t-1,t+1]$ contains $O(\log(q(|t|+2)))$ nontrivial zero ordinates counted with multiplicity.
\end{lemma}
\begin{proof}
Let $a\in\{0,1\}$ be the parity of $\chi$ and set
\[
\Xi_\chi(s)=(s(s-1))^{\delta_\chi}
(q/\pi)^{(s+a)/2}\Gamma((s+a)/2)L(s,\chi).
\]
This is entire of order one and its zeros are exactly the nontrivial zeros. The functional equation and conjugation give $\Xi_\chi(s)=\varepsilon_\chi\overline{\Xi_\chi(1-\bar s)}$ with a constant of modulus one. For $s_0=1/2+it$ not a zero, differentiation implies $\Rea(\Xi_\chi'/\Xi_\chi)(s_0)=0$. Hadamard's product gives the absolutely convergent difference identity
\begin{equation}\label{eq:had}
\frac{\Xi_\chi'}{\Xi_\chi}(s)-\frac{\Xi_\chi'}{\Xi_\chi}(2+it)
=\sum_\rho\left(\frac1{s-\rho}-\frac1{2+it-\rho}\right)
\end{equation}
for $s$ away from zeros. In the difference, the linear exponential and canonical-product regularizers cancel. Absolute convergence follows from the genus-one product (or its order-one zero count). On GRH, every $1/(s_0-\rho)$ is purely imaginary. Thus at any evaluation point with $\Xi_\chi(s_0)\ne0$,
\begin{equation}\label{eq:local}
\Rea\frac{\Xi_\chi'}{\Xi_\chi}(2+it)
=\sum_\rho\frac{3/2}{(3/2)^2+(t-\Im\rho)^2}.
\end{equation}
At $2+it$, the Euler series bounds $L'/L$ absolutely by $\sum_n\Lambda(n)n^{-2}$. The gamma logarithmic derivative is $O(\log(|t|+2))$ and the conductor term is $\frac12\log(q/\pi)$; the two rational terms for $\zeta$ are uniformly bounded. Hence the left side of \eqref{eq:local} is $O(\log(q(|t|+2)))$. For an arbitrary $t$, choose $t'$ within $1/2$ of $t$ that is not a zero ordinate for this character. Each zero with $|\Im\rho-t|\le1$ contributes a fixed positive lower bound at $t'$. This proves the claimed local count with an absolute constant, including the central zero.

We next derive the contour bounds. For $s=1/4+it$, the summands of \eqref{eq:had} at $|t-\Im\rho|\le1$ are $O(1)$, using the real separation $1/4$ from the critical line. For the rest, the difference of reciprocals is $O(|t-\Im\rho|^{-2})$. The local count and summation over unit intervals bound this tail by $O(\log(q(|t|+2)))$: one uses $\log(q(|t|+k+2))\le\log(q(|t|+2))+\log(k+1)$ for integers $k\ge1$. Returning to $L'/L$ by subtracting the gamma and elementary terms gives
\[
|L'/L(1/4+it,\chi)|\ll\log(q(|t|+2)).
\]
For the horizontal edges, for each integer $M\ge3$ choose heights $t_M^+\in[M,M+1]$ and $t_M^-\in[-M-1,-M]$ separated from all zero ordinates by $\gg1/\log(q(M+2))$. These heights exist by deleting small intervals about the $O(\log(q(M+2)))$ ordinates within distance one of the height interval. On $1/4\le\Rea s\le2$ at these heights, the same difference argument gives $L'/L(s,\chi)\ll\log^2(q(M+2))$. The zeta pole is away from the edges.

For $y>0$ the inverse Mellin integral on $\Rea s=2$ of $y^s/(s(s+1))$ equals $(1-y^{-1})_+$, including zero at $y=1$. This follows by closing the contour to the left for $y>1$ and to the right for $y<1$; at $y=1$ use continuity of the absolutely convergent integral. The Euler product and absolute convergence permit termwise integration, so
\[
\sum_{n\le x}\Lambda(n)\chi(n)(1-n/x)
=\frac1{2\pi i}\int_{(2)}-\frac{L'}{L}(s,\chi)\frac{x^s}{s(s+1)}\,ds.
\]
Shift to $\Rea s=1/4$ along the selected heights. For fixed $q,x$ the horizontal integrals tend to zero, bounded by a constant times $x^2\log^2(q(M+2))/M^2$. The crossed poles are the principal pole at $1$ (residue $x/2$ for the integrand) and the nontrivial zeros (residues $-m x^\rho/[\rho(\rho+1)]$ for zeros of order $m$). The kernel poles at $0,-1$ and trivial zeros are outside the contour. The new vertical integral has absolute value at most a constant times
\[
x^{1/4}\int_{-\infty}^{\infty}\frac{\log(q(|t|+2))}{1+t^2}\,dt
\ll x^{1/4}\log(q+2).
\]
The local count and the critical-line identity $|\rho(\rho+1)|\asymp1+|\Im\rho|^2$ imply both bounds of \eqref{eq:zerotail} by summing the unit intervals; for the tail use the integral comparison for $(\log(q(u+2)))/(1+u^2)$. This proves absolute convergence of the residue sum and completes the contour argument.
\end{proof}

\section{Proof of the universal reconstruction theorem}
The coefficients from Section~\ref{sec:code} satisfy $\sum_j c_j\log(q_j+2)<\infty$, since $\log(q_j+2)\le q_j+1$. The bounds on $W$ follow from $\chi_0(n)=1$, $|\chi_j(n)|\le1$ and \eqref{eq:total}. At every $x$, absolute convergence of the weight series allows it to be interchanged with the finite sum over $n$. Apply Lemma~\ref{lem:explicit} separately to each primitive character and sum, using the absolute uniform bounds. On setting $x=\e^t$, this yields for $t\ge0$
\begin{equation}\label{eq:almost}
F(t)=G(t)+O(\e^{-t/4}),\qquad
G(t)=\sum_{\gamma\in\mathcal Z} a_\gamma \e^{i\gamma t},\quad
a_\gamma=-\frac{A_\gamma}{\rho_\gamma(\rho_\gamma+1)},\quad
A_\gamma=\sum_j c_jm_j(\gamma),
\end{equation}
where $\mathcal Z$ denotes the countable union of distinct real ordinates over all primitive characters. The $O$-constant is finite and independent of $t$. For all $Y\ge1$ we have
\begin{equation}\label{eq:sumtail}
B:=\sum_{\gamma\in\mathcal Z}|a_\gamma|<\infty,\qquad
\sum_{|\gamma|>Y}|a_\gamma|\ll\frac{\log(Y+2)}{Y+1}.
\end{equation}
Indeed all $c_j$ and $m_j(\gamma)$ are nonnegative, so the sums of absolute values regroup without any cancellation or conditional-convergence issue. Bound them using \eqref{eq:zerotail}, and $\sum_jc_j\log(q_j+2)<\infty$. A single ordinate may appear for infinitely many characters, and there may be infinitely many distinct ordinates in a bounded interval across all conductors. Neither issue affects absolute convergence or the following argument. In particular $G$ is uniformly bounded and its series converges uniformly on $\R$.

For each fixed $h$, replacing $F(t+h)\overline{F(t)}$ by $G(t+h)\overline{G(t)}$ changes its mean over $[0,T]$ by $o(1)$: both $G$ and $F$ on $[0,\infty)$ are bounded, and the difference in \eqref{eq:almost} decays exponentially. A fixed initial interval where $t+h<0$ is harmless, since the functions are locally bounded. Expand the product using the absolutely summable double series, of total absolute mass at most $B^2$. Each term contains
\[
\frac1T\int_0^T\e^{i(\gamma-\eta)t}\,dt\longrightarrow
\begin{cases}1&\gamma=\eta,\\0&\gamma\ne\eta.\end{cases}
\]
These integrals have modulus at most one. Dominated convergence for the countable double sum, or uniform approximation by finite sums, yields
\begin{equation}\label{eq:correlation}
C(h)=\sum_{\gamma\in\mathcal Z}|a_\gamma|^2\e^{i\gamma h}.
\end{equation}
The same argument for \eqref{eq:p}, using $\sum|a_\gamma|^2\le B^2$, proves $P(\gamma)=|a_\gamma|^2$ for $\gamma\in\mathcal Z$, and $P(\gamma)=0$ otherwise. Thus $P$ is real and nonnegative, and positivity of $A_\gamma$ gives \eqref{eq:recover} even when $\gamma=0$.

To apply Proposition~\ref{prop:code}, use the local count in Lemma~\ref{lem:explicit}:
\[
m_j(\gamma)\le C_1\log(q_j(|\gamma|+2))
\le H_\gamma (q_j+1),
\]
where one can choose an integer $H_\gamma\ge C_1(1+\log(|\gamma|+2))$. The constant $C_1$ is absolute. Therefore the multiplicity sequence lies in $\mathcal M$ at every ordinate. It is uniquely recoverable by the code, even if infinitely many of its entries are nonzero. The set of ordinates is the support in the sense of nonzero atoms of $P$, not necessarily a topologically discrete set. GRH locates all nontrivial zeros at the recovered ordinates on the critical line. Finally, for any imprimitive character the extra finite Euler factors have zeros only on $\Rea s=0$; hence the critical-strip zero multiset is inherited. This proves Theorem~\ref{thm:main}.

\section{Only ordinary primes are needed}
Define $S_{\rm p}(x)=\sum_{p\le x} (\log p) W(p)(1-p/x)$ and the known constant
\[
b=\frac23\sum_{\{j:\chi_j^2\text{ is principal modulo }q_j\}}c_j.
\]
For $q_j=1$ the condition is understood to hold. The series defining $b$ is convergent and computably approximable by finite character data.
\begin{corollary}[Prime-only reconstruction]\label{cor:prime}
Theorem~\ref{thm:main} is unchanged if $F$ is replaced by
\[
F_{\rm p}(t)=\e^{-t/2}\bigl(S_{\rm p}(\e^t)-\e^t/2\bigr)+b.
\]
In particular the two signals have exactly the same limiting autocorrelation.
\end{corollary}
\begin{proof}
For each fixed primitive $\chi_j$, the classical prime number theorem in arithmetic progressions gives
\[
\theta_j^{(2)}(y):=\sum_{p\le y}\chi_j(p)^2\log p
=\epsilon_j y+o_j(y),
\]
where $\epsilon_j$ is one if the squared character is principal and zero otherwise. This follows by character orthogonality from the prime number theorem in reduced residue classes. With $y=\sqrt{x}$, partial summation gives
\begin{align*}
\sum_{p\le y}\chi_j(p)^2\log p\,(1-p^2/y^2)
&=\frac2{y^2}\int_1^y u\theta_j^{(2)}(u)\,du\\
&=\tfrac23\epsilon_j y+o_j(y).
\end{align*}
Chebyshev's bound $\theta(y)\ll y$ implies that the normalized expression on the left divided by $y$ is absolutely bounded, uniformly in $j$ and $y\ge1$. Dominated convergence with $\sum_j c_j<\infty$ therefore gives the aggregate prime-square contribution $b\sqrt{x}+o(\sqrt{x})$. For powers of exponent $k\ge3$, there are at most $\log x/\log 2$ possible $k$, and $\sum_{p\le x^{1/k}}\log p\le x^{1/3}\log x$. Since $|W|\le5/4$, their contribution is $O(x^{1/3}\log^2 x)=o(\sqrt{x})$. The von Mangoldt weight at each prime power is $\log p$, as used in both estimates.

Consequently $F_{\rm p}(t)-F(t)=o(1)$ as $t\to\infty$. Both are locally bounded and bounded for $t\ge0$. Their autocorrelation difference has Ces\`aro mean zero at every fixed lag, proving the corollary.
\end{proof}
\begin{remark}
Even without adding $b$ to the normalized prime-only signal one can decode: all nonzero frequencies are unchanged, and its central Fourier coefficient is $-\tfrac43 A_0-b$. Since $A_0,b\ge0$, the square root of its zero-frequency power determines $A_0$ by subtracting $b$ and multiplying by $3/4$. The explicit correction in Corollary~\ref{cor:prime} is convenient because the formula is identical at all frequencies.
\end{remark}

\section{A finite-family converse}
The countable theorem is a consequence of GRH. A related finite-character probe also detects failure of GRH, so this smoothing has a zero-location interpretation analogous to the reduction in the Lean proof.
\begin{proposition}\label{prop:converse}
Let $J$ be a finite set of primitive characters, with strictly positive real weights $c_j$, and include the conductor-one character if desired. Let $S_J(x)$ be the corresponding finite weighted prime-power sum with the same kernel, and let $d$ be the coefficient of the conductor-one character (zero if absent). Then
\[
S_J(x)-dx/2=O(\sqrt{x})
\]
holds if and only if all the nontrivial zeros of the characters in $J$ are on the critical line.
\end{proposition}
\begin{proof}
The forward implication from GRH to the bound is immediate from Lemma~\ref{lem:explicit}. Conversely, for $\Rea s>1$ termwise integration (absolute convergence) gives
\[
\int_1^\infty (S_J(x)-dx/2)x^{-s-1}\,dx
=\frac{\sum_{j\in J}c_j(-L'/L)(s,\chi_j)}{s(s+1)}-\frac{d}{2(s-1)}.
\]
The assumed bound extends the integral holomorphically to $\Rea s>1/2$ by locally uniform absolute convergence. The right-hand side is meromorphic there; the pole at $1$ cancels by the choice of main term and standard nonvanishing on $\Rea s=1$. At any nontrivial zero $\rho$ to the right of the line, the residue of the right-hand side is
\[
-\frac{\sum_{j\in J}c_j\ord_\rho L(s,\chi_j)}{\rho(\rho+1)}\ne0.
\]
No cancellation is possible because the numerator is positive. This contradicts holomorphy. The functional equation and conjugation reflect each character's zeros about the critical line, so the left half of the critical strip is also excluded. This proves the equivalence.
\end{proof}
No analogous converse for the countable aggregate is asserted here: meromorphic continuation of a countable sum through possible accumulation of zeros is not automatic. Nor is a one-sided bound without an absolute-value or boundedness hypothesis asserted in this proposition.

\section{Stability, numerical limits, and further questions}
For an observed ordinate $\gamma$, put $K_\gamma=|\rho_\gamma(\rho_\gamma+1)|$. If a nonnegative approximation $P'$ satisfies
\[
|P'-P(\gamma)|<\frac{d_n^2}{64K_\gamma^2},
\]
then $|K_\gamma\sqrt{P'}-A_\gamma|<d_n/8$ because $|\sqrt{u}-\sqrt{v}|\le\sqrt{|u-v|}$ for $u,v\ge0$. Thus for $n+1\ge H_\gamma$ Proposition~\ref{prop:code} gives a finite prefix decoding guarantee. The code is not designed for efficient numerical evaluation: its precision thresholds can be extremely small. This guarantee concerns errors in the spectral power at an exact specified frequency, not errors in the frequency itself.

The existence of both limiting averages is independent of any gap assumptions. For finite-window analysis one may truncate by \emph{character index and height} to a finite set $E$ of frequencies, retaining the full aggregate coefficients at those frequencies in the theoretical analysis. For $\epsilon=\sum_{\gamma\notin E}|a_\gamma|$, $B$ as in \eqref{eq:sumtail}, and $\Delta=\min_{\gamma\ne\eta\in E}|\gamma-\eta|>0$ (omit the spacing term if $|E|\le1$), one has for $h\ge0$
\[
\left|\frac1T\int_0^T F(t+h)\overline{F(t)}\,dt-C(h)\right|
\le O(1/T)+3B\epsilon+\frac{2B^2}{T\Delta}.
\]
Indeed the exponentially decaying remainder contributes $O(1/T)$ uniformly in such $h$. Truncating the double product costs at most $2B\epsilon$, the omitted diagonal costs at most $B\epsilon$, and each off-diagonal exponential average is at most $2/(T\Delta)$. Such finite sets exist by absolute convergence; conductor and height truncation have vanishing tails by \eqref{eq:zerotail} and \eqref{eq:total}. This estimate does not give a useful uniform complexity bound without control of frequency separations. The exact infinite-statistic theorem is not a finite-data algorithm for discovering arbitrary zeros.

Further research questions:
\begin{enumerate}[leftmargin=*]
\item Can the coefficient sequence be optimized to improve the prefix precision bound while still distinguishing multiplicities across unbounded conductors?
\item Can one encode the same labelled spectrum with a single \emph{real} nonnegative prime weight, perhaps retaining higher-order correlations to remove the conjugation ambiguity of real autocorrelation?
\item Can one give effective finite-window reconstruction of a prescribed finite collection of characters without zero-separation assumptions, using interval certificates and combined conductor/height truncations?
\item Can the finite-family zero-detection converse be extended to a countable aggregate, or a one-sided estimate in the style of the original proof, without assuming unjustified meromorphic continuation?
\item Does the universal encoding extend to suitable families of automorphic $L$-functions with uniform zero counts and explicit formulas, conditional on the corresponding zero-location theorem? That GRH is not supplied by the Dirichlet proof.
\item Can the finite coding lemma, dominated-convergence argument and prime-only correction be formally verified in Lean as modular extensions of the exported statements?
\end{enumerate}

\section*{Verification and novelty statement}
The accompanying Python script checks the rational coefficient recursion, exact interval bounds on prefix sums, and recovery of synthetic multiplicities including a shared ordinate. It also checks the Fourier averaging identity numerically for a finite synthetic spectrum. These tests are reproducible sanity checks, not proof of the analytic theorems or a computation of actual $L$-function zeros. The full proof is in Sections 3--7; Section 8 adds stability and research directions. The novel candidate is the universal character-labelled inverse result with countable positive encoding and the ordinary-prime correction; the explicit formula, Hadamard factorization, algebraic norm bound, and Bohr orthogonality are established mathematical techniques. A limited search for related prime-race, autocorrelation and explicit-formula work is not a comprehensive priority search. No Lean build was performed in this environment.

\begin{thebibliography}{9}
\bibitem{repo} L. Sernau et al., \emph{GRH for Dirichlet L-functions, in Lean 4}, 25 November 2026. Commit \texttt{72fead102f663d5c7a407964a13dc0cf9d977207}. \url{https://github.com/lsernau/grh-lean}. See \texttt{grh/GUIDE.md}, \texttt{grh/src/Statements.lean}, \texttt{grh/src/Main.lean}, and \texttt{grh/src/New/Reduction/} files \texttt{MainInequality.lean} and \texttt{OneSidedBound.lean}.
\bibitem{wiki} Wikipedia, \emph{Riemann hypothesis}, accessed 26 November 2026. \url{https://en.wikipedia.org/wiki/Riemann_hypothesis}.
\bibitem{rub} M. Rubinstein, \emph{Alp\"oge/Claude proof of the Generalized Riemann Hypothesis}, independent verification records, November 2026. \url{https://www.math.uwaterloo.ca/~mrubinst/alpoge_claude_grh/}.
\bibitem{mv} H. L. Montgomery and R. C. Vaughan, \emph{Multiplicative Number Theory I: Classical Theory}, Cambridge University Press, 2007, chapters on analytic properties of Dirichlet $L$-functions, explicit formulas and primes in progressions.
\bibitem{dav} H. Davenport, \emph{Multiplicative Number Theory}, 3rd ed., revised by H. L. Montgomery, Springer, 2000.
\bibitem{ans} A. Akbary, N. Ng and M. Shahabi, Limiting distributions of the classical error terms of prime number theory, \emph{Q. J. Math.} \textbf{65} (2014), 743--780. \url{https://arxiv.org/abs/1306.1657}.
\bibitem{rs} M. Rubinstein and P. Sarnak, Chebyshev's bias, \emph{Experimental Mathematics} \textbf{3} (1994), 173--197.
\end{thebibliography}
\end{document}
